% Quantitative diagrams: shared formulas determine points and curves.
\newcommand{\vdWIsotherms}{%
\begin{center}\begin{tikzpicture}
\begin{axis}[width=.82\linewidth,height=6cm,axis lines=left,xmin=.48,xmax=3.5,ymin=0,ymax=2.7,
xlabel={$v'=V/V_c$},ylabel={$p'=p/p_c$},xtick={1,2,3},ytick={1,2},clip=true]
\foreach \t in {.85,.9,1,1.08,1.16,1.24}{
\addplot[DeepBlue,domain=.48:3.3,samples=250]{8*\t/(3*x-1)-3/x^2};}
\addplot[only marks,mark=*,mark size=1.3pt] coordinates{(1,1)};
\node[above right] at (axis cs:1.02,1.04) {$C$};
\draw[dashed,TextGray] (axis cs:1,0)--(axis cs:1,1);
\node[align=left,anchor=north east,font=\small] at (axis cs:3.35,2.5)
{$t'=0.85,\ 0.90,\ 1.00,$\\$1.08,\ 1.16,\ 1.24$\\从下向上};
\end{axis}\end{tikzpicture}\end{center}}

\newcommand{\criticalPointFigure}{%
\begin{center}\begin{tikzpicture}
\begin{axis}[width=6cm,height=5.3cm,axis lines=left,xmin=.55,xmax=2.3,ymin=0,ymax=2.4,
xlabel={$v'$},ylabel={$p'$},xtick={1,2},ytick={1,2}]
\addplot[DeepBlue,thick,domain=.56:2.2,samples=250]{8/(3*x-1)-3/x^2};
\addplot[only marks,mark=*] coordinates{(1,1)};
\draw[dashed,TextGray] (axis cs:.8,1)--(axis cs:1.3,1);
\node[above right] at (axis cs:1.07,1.1) {$C$};
\end{axis}
\begin{scope}[xshift=7cm,x=.88cm,y=.82cm]
\draw[axis] (0,0)--(5.7,0) node[right]{$T$};
\draw[axis] (0,0)--(0,5.5) node[above]{$p$};
\draw[DeepBlue,thick,domain=.6:4,samples=100] plot(\x,{.2+.2*\x*\x});
\coordinate (P) at (1,1.5);\coordinate (M) at (2.54951,1.5);\coordinate (Q) at (5,1.5);
\coordinate (A) at (1,4.6);\coordinate (B) at (5,4.6);\coordinate (C) at (4,3.4);
\draw[->,Ink] (P)--(M);\draw[->,Ink] (M)--(Q);
\draw[->,DeepBlue,dashed] (P)--(A);\draw[->,DeepBlue,dashed] (A)--(B);\draw[->,DeepBlue,dashed] (B)--(Q);
\foreach \q in {P,M,Q,A,B,C}{\fill (\q) circle(1.2pt);}
\foreach \q in {P,M,Q}{\node[below=4pt] at (\q){$\q$};}
\foreach \q in {A,B,C}{\node[above=4pt] at (\q){$\q$};}
\node[font=\small] at (2.2,3.7){液相};\node[font=\small] at (4.7,.55){气相};
\end{scope}
\node[font=\small] at (2.5,-1) {(a) 临界等温线};
\node[font=\small] at (9,-1) {(b) 穿越与绕行路径};
\end{tikzpicture}\end{center}}

\newcommand{\maxwellAreaFigure}{%
\begin{center}\begin{tikzpicture}
\begin{axis}[width=.85\linewidth,height=6cm,axis lines=left,xmin=.48,xmax=2.9,ymin=.25,ymax=1.08,
xlabel={$v'$},ylabel={$p'$},ylabel style={at={(axis description cs:0,1)},anchor=south east},xtick={.603402,2.348842},xticklabels={$v_l'$,$v_g'$},ytick={.64699835},yticklabels={$p_T'$}]
\addplot[draw=none,fill=DeepBlue!15,domain=.603401903:1.090526633,samples=120]
{7.2/(3*x-1)-3/x^2} -- (axis cs:1.090526633,.646998352)--(axis cs:.603401903,.646998352)\closedcycle;
\addplot[draw=none,fill=Gold!20,domain=1.090526633:2.348842376,samples=120]
{7.2/(3*x-1)-3/x^2} -- (axis cs:2.348842376,.646998352)--(axis cs:1.090526633,.646998352)\closedcycle;
\addplot[DeepBlue,thick,domain=.54:2.8,samples=300]{7.2/(3*x-1)-3/x^2};
\draw[Ink,thick] (axis cs:.603401903,.646998352)--(axis cs:2.348842376,.646998352);
\addplot[only marks,mark=*,mark size=1.2pt] coordinates{(.603402,.646998)(.718597,.419843)(1.090527,.646998)(1.528505,.724013)(2.348842,.646998)};
\node[above left] at (axis cs:.603402,.646998){$M$};
\node[below] at (axis cs:.718597,.419843){$A$};
\node[above left] at (axis cs:1.090527,.646998){$P$};
\node[above] at (axis cs:1.528505,.724013){$B$};
\node[above right] at (axis cs:2.348842,.646998){$N$};
\node at (axis cs:.78,.58){$S_1$};\node at (axis cs:2.05,.83){$S_2$};
\draw[->,TextGray] (axis cs:2.05,.79)--(axis cs:1.85,.681);
\node[font=\small] at (axis cs:2.2,.96){$t'=0.9,\ S_1=S_2$};
\end{axis}\end{tikzpicture}\end{center}}

\newcommand{\contrailFigure}{%
\begin{center}\includegraphics[width=.88\linewidth]{/Users/dongxingxi/Documents/物理竞赛一轮资料/sample/xes-original/restructured-complete/assets/ch17/contrail.png}\end{center}}

\newcommand{\mercuryTubeFigure}{%
\begin{center}\begin{tikzpicture}[x=.17cm,y=.17cm]
\draw[Ink,thick] (0,5)--(0,10)--(31,10) (0,5)--(31,5);
\fill[DeepBlue!25] (17,5.1) rectangle(27,9.9);
\draw[DeepBlue] (17,5)--(17,10) (27,5)--(27,10);
\draw[<->] (0,2)--(17,2) node[midway,below,font=\small]{$17\,\mathrm{cm}$};
\draw[<->] (17,13)--(27,13) node[midway,above,font=\small]{$10\,\mathrm{cm}$};
\node[font=\small] at (22,7.5){Hg};
\node[font=\small,align=center] at (8.5,18){封闭气体\\空气与饱和汞蒸气};
\draw[->,TextGray] (8.5,15)--(8.5,9);
\draw[Ink,thick] (47,0)--(47,33)--(52,33)--(52,0);
\fill[DeepBlue!25] (47.1,3) rectangle(51.9,13);
\draw[DeepBlue] (47,3)--(52,3) (47,13)--(52,13);
\node[font=\small] at (49.5,8){Hg};
\draw[<->] (56,13)--(56,33) node[midway,right,font=\small]{$20\,\mathrm{cm}$};
\draw[<->] (56,3)--(56,13) node[midway,right,font=\small]{$10\,\mathrm{cm}$};
\node[font=\small,align=center] at (37,26){同温度的\\封闭气体};
\draw[->,TextGray] (44,24)--(48,24);
\node[below,font=\small] at (49.5,0){开口，外接大气};
\node[above,font=\small] at (49.5,34){封闭端};
\end{tikzpicture}\end{center}}

\newcommand{\mixedGasPVFigure}{%
\begin{center}\begin{tikzpicture}
\begin{axis}[width=.8\linewidth,height=5.7cm,axis lines=left,xmin=0,xmax=4.6,ymin=0,ymax=5.7,
xlabel={$x=V'/V$},ylabel={$p/\mathrm{atm}$},xtick={.5,1,2,4},ytick={1,2,3,5}]
\addplot[DeepBlue,thick,domain=.45:2,samples=150]{1+2/x};
\addplot[DeepBlue,thick,dashed,domain=2:4.4,samples=120]{4/x};
\addplot[only marks,mark=*] coordinates{(.5,5)(1,3)(2,2)(4,1)};
\draw[dotted] (axis cs:2,0)--(axis cs:2,2);
\node[align=left,font=\small] at (axis cs:3.2,4.3){$x<2$：液体仍在\\$x\ge2$：水全汽化};
\node[above right,font=\small] at (axis cs:2,2.1){液体刚好消失};
\end{axis}\end{tikzpicture}\end{center}}

\newcommand{\pressureCookerQuestionFigure}{%
\begin{center}\begin{tikzpicture}
\begin{scope}[x=.6cm,y=.6cm]
\fill[DeepBlue!15] (.05,.05) rectangle(3.95,1.7);
\draw[thick] (0,4)--(0,0)--(4,0)--(4,4);
\draw[thick] (-.2,4)--(1.7,4)--(1.7,4.6) (2.3,4.6)--(2.3,4)--(4.2,4);
\draw[thick,fill=gray!15] (1.4,4.65) rectangle(2.6,5.15);
\draw[DeepBlue] (0,1.7)--(4,1.7);
\node[font=\small,above] at (2,5.3){限压阀 $S$};
\node[font=\small] at (2,2.6){气相};
\node[font=\small] at (2,.8){液态水};
\end{scope}
\begin{axis}[at={(3.6cm,0)},anchor=south west,width=8cm,height=5.0cm,axis lines=left,
xmin=0,xmax=130,ymin=0,ymax=260,xlabel={$t/{}^\circ\mathrm C$},ylabel={$p_w/\mathrm{kPa}$},
xtick={10,30,50,70,90,100,120},ytick={0,50,100,150,200,250},grid=both,
minor tick num=1,grid style={draw=RuleGray!70},tick label style={font=\scriptsize}]
\addplot[DeepBlue,thick] table[x=t,y=ps]{assets/ch17/saturation.dat};
\end{axis}\end{tikzpicture}\end{center}}

\newcommand{\pressureCookerFigure}{%
\begin{center}\begin{tikzpicture}
\begin{scope}[x=.7cm,y=.7cm]
\fill[DeepBlue!15] (.05,.05) rectangle(3.95,1.7);
\draw[thick] (0,4)--(0,0)--(4,0)--(4,4);
\draw[thick] (-.2,4)--(1.7,4)--(1.7,4.6) (2.3,4.6)--(2.3,4)--(4.2,4);
\draw[thick,fill=gray!15] (1.4,4.65) rectangle(2.6,5.15);
\draw[DeepBlue] (0,1.7)--(4,1.7);
\draw[->,DeepBlue] (2,3.2)--(2,4.5);
\draw[->,DeepBlue] (1.7,4.63)--(.7,4.63);
\node[font=\small,above] at (2,5.3){限压阀 $S$};
\node[font=\small,align=center] at (2,2.4){空气与水蒸气};
\node[font=\small] at (2,.8){液态水};
\end{scope}
\begin{axis}[at={(4.2cm,0)},anchor=south west,width=8cm,height=5.8cm,axis lines=left,
xmin=0,xmax=130,ymin=0,ymax=340,xlabel={$t/{}^\circ\mathrm C$},ylabel={$p/\mathrm{kPa}$},
xtick={10,90,120},ytick={100,200,300},legend style={at={(.02,.98)},anchor=north west,font=\scriptsize,draw=none,fill=none}]
\addplot[DeepBlue,thick] table[x=t,y=ps]{assets/ch17/saturation.dat};\addlegendentry{$p_w$}
\addplot[TextGray,dashed] table[x=t,y=pa]{assets/ch17/saturation.dat};\addlegendentry{$p_a$（封闭外推）}
\addplot[Ink,thick,restrict x to domain=0:90.174448744] table[x=t,y=total]{assets/ch17/saturation.dat};\addlegendentry{$p_a+p_w$（开启前）}
\addplot[Ink,dashed,restrict x to domain=90.174448744:130] table[x=t,y=total]{assets/ch17/saturation.dat};
\addplot[Ink,thick,domain=90.174448744:120]{198.671424787};
\addplot[dotted,domain=0:128]{198.671424787};
\addplot[only marks,mark=*] coordinates{(90.174448744,198.671424787)(120,198.671424787)};
\draw[dotted] (axis cs:90.174448744,0)--(axis cs:90.174448744,198.671424787)
(axis cs:120,0)--(axis cs:120,198.671424787);
\node[font=\small,anchor=north] at (axis cs:109,191){$p_*$};
\end{axis}\end{tikzpicture}\end{center}}

\newcommand{\twoVesselsFigure}{%
\begin{center}\begin{tikzpicture}[x=1cm,y=1cm]
\fill[DeepBlue!12] (7.22,.02) rectangle(10.38,.95);
\draw[thick] (3.2,1.5)--(3.2,0)--(0,0)--(0,3.45)--(3.2,3.45)--(3.2,1.9);
\draw[thick] (7.2,1.5)--(7.2,0)--(10.4,0)--(10.4,3.45)--(7.2,3.45)--(7.2,1.9);
\draw[thick] (3.2,1.5)--(7.2,1.5) (3.2,1.9)--(7.2,1.9);
\draw[DeepBlue] (7.2,.95)--(10.4,.95);
\node[font=\small,above] at (5.2,2.15){连通细管};
\node[align=center] at (1.6,2.6){热端\\$100{}^\circ\mathrm C$};
\node[align=center] at (8.8,2.6){冷端\\$0{}^\circ\mathrm C$};
\node[font=\small] at (8.8,.45){凝结水};
\draw[->,DeepBlue] (4.2,.9)--(6.2,.9);
\node[font=\small,below] at (5.2,.8){水蒸气向冷端迁移};
\draw[TextGray] (7,-.2)--(10.6,-.2)--(10.6,1.2) (7,-.2)--(7,1.2);
\node[font=\small,below] at (8.8,-.25){外部冰水浴};
\end{tikzpicture}\end{center}}

\newcommand{\cloudFormationFigure}{%
\begin{center}\begin{tikzpicture}
\begin{scope}[x=.8cm,y=.8cm]
\draw[axis] (0,0)--(0,5.5) node[above]{$z$};
\foreach \y/\r in {.8/.28,2/.36,3.2/.44,4.4/.52}{\draw[DeepBlue] (1.5,\y) circle(\r);}
\draw[->,DeepBlue] (2.4,.7)--(2.4,4.6);
\draw[dashed,TextGray] (0,4.4)--(3.8,4.4);
\node[font=\small,align=center] at (1.8,5.5){云底 $z_d$\\首次饱和};
\node[font=\small,align=center] at (1.8,-.5){上升膨胀、降温};
\end{scope}
\begin{axis}[at={(4.2cm,0)},anchor=south west,width=8cm,height=5.5cm,axis lines=left,
xmin=269,xmax=290,ymin=.25,ymax=1.12,xlabel={$T/\mathrm K$},ylabel={$p_v/p_s(T_0)$},
xtick={270,278.475,288},xticklabels={270,$T_d$,288},ytick={.4,.6,.8,1}]
\addplot[DeepBlue,thick,domain=270:288,samples=200]{exp(44000/8.314*(1/288-1/x))};
\addplot[Ink,thick,domain=278.474927:288,samples=120]{.6*(x/288)^3.5};
\addplot[Ink,dashed,domain=270:278.474927,samples=80]{.6*(x/288)^3.5};
\addplot[only marks,mark=*] coordinates{(278.474927,.533370375)};
\draw[dotted] (axis cs:278.474927,.25)--(axis cs:278.474927,.533370375);
\node[font=\small] at (axis cs:280,.99){$p_s(T)$};
\node[font=\small] at (axis cs:285,.44){$p_v(T)$};
\draw[->,TextGray] (axis cs:287,.67)--(axis cs:282,.62);
\end{axis}\end{tikzpicture}\end{center}}
